\documentclass[11pt]{article}
\usepackage[margin=1in]{geometry}
\usepackage{amsmath,amssymb,amsthm}
\usepackage{hyperref}

\title{Partial Results in Empirical Complexity:\\
\large Template for PNP Research Laboratory Reports}
\author{PNP Research Laboratory}
\date{\today}

\newtheorem{conjecture}{Conjecture}
\newtheorem{lemma}{Lemma}
\newtheorem{observation}{Observation}

\begin{document}
\maketitle

\begin{abstract}
This template separates \emph{observations}, \emph{conjectures}, and \emph{theorems}.
Experimental claims are at most evidence level E2 and do not resolve $\mathrm{P}$ vs $\mathrm{NP}$.
\end{abstract}

\section{Evidence statement}
Every numbered claim below is tagged with an evidence level (E0--E6)
as defined in the laboratory methodology.

\section{Definitions}
Cite Arora--Barak / Sipser; do not invent definitions.

\section{Experimental setup}
Record seeds, generator, $n$, ratios, solvers, timeouts, code commit.

\begin{observation}[E2]
State only finite-sample measurements here.
\end{observation}

\begin{conjecture}[E3]
State a precise mathematical conjecture. Prioritize attempted refutation.
\end{conjecture}

\section{Why this does not settle P vs NP}
Explicit barriers: finite vs asymptotic; average vs worst case; \ldots

\section{Reproducibility}
Exact CLI and artifact hashes.

\section*{Disclaimer}
AI-assisted text must be labelled and independently audited.
\end{document}
